\chapter{Gérer un livre exotique}\label{ch:dv:managing-an-exotic-book}

Un desk vend 100 millions d'un \omterm{def:dv:autocallables:autocallable}{autocallable} worst-of à trois ans à ses clients de banque privée, au pair. Son modèle valorise les titres à 98, de sorte que la vente apparaisse
comme un profit de 2 millions. Le contrôle de valorisation en reconnaît 325\,000. Il passe d'un coup 621\,000 pour le coût de clôture des couvertures, et retient
1\,054\,000 jusqu'à ce que deux éléments qu'il ne peut pas observer deviennent observables~: quel modèle de volatilité convient à ce payoff, et la corrélation entre les
deux indices. Un an plus tard, la corrélation se traite, et 424\,000 sont repris au résultat. Ce chapitre porte sur cette arithmétique et sur le livre qui la sous-tend~: ce que détient un livre exotique,
les réserves sur sa valorisation, le risque de modèle tel que le desk le rencontre, le stress et la concentration, ainsi que le \omterm{def:dv:managing-an-exotic-book:dayone}{P\&L du premier jour} et sa reprise.

\section{Ce que détient un livre exotique}

Un livre exotique est une vaste position vendeuse sur des \omterm{def:dv:the-structured-products-business:product}{produits structurés} (chapitres 18 et 19) et sur les couvertures achetées en contrepartie, vanilles, variance et
le sous-jacent. Les \omterm{def:dv:greeks-and-the-hedging-pnl:greeks}{grecques} nettes sont faibles~; ce qui subsiste, ce sont les expositions que personne ne vend. Pour les \omterm{def:dv:autocallables:autocallable}{autocallables} vendus aux investisseurs, il s'agit d'une position longue en volatilité aux
échéances longues, position longue sur le skew et position vendeuse en corrélation sur les worst-of, le tout concentré là où se trouvent les barrières des titres.

La comptabilité classe ces expositions selon les données d'entrée dont leur valorisation a besoin. La norme IFRS 13 hiérarchise ces données en trois niveaux~: les cours cotés d'instruments identiques sur
des marchés actifs (niveau 1)~; les autres données observables directement ou indirectement, pour la quasi-totalité de la durée de l'instrument (niveau 2)~; les données non observables
(niveau 3). Une corrélation à trois ans entre deux indices n'a pas de marché~: le worst-of est un instrument de niveau 3, et les règles sur son profit du premier jour
découlent de ce fait.

\begin{example}[Le titre]\label{ex:dv:managing-an-exotic-book:note}
\omterm{def:dv:autocallables:autocallable}{Autocallable} Phoenix worst-of sur deux indices (volatilités 20~\% et 25~\%, corrélation 0.5, taux 3~\%, dividendes 2~\%), trois dates d'observation annuelles,
rappel automatique à 100~\%, un coupon de 7.95~\% par an versé si l'indice est au moins égal à 70~\% avec mémoire, capital protégé à l'échéance sauf si le moins bon indice est inférieur à 60~\%.
Pour ce coupon, le modèle valorise le titre à 98.00 pour 100. Du côté du desk, qui est vendeur du titre, le livre est long de 0.533 et 0.693 de \omterm{def:dv:greeks-and-the-hedging-pnl:greeks}{vega} par
point de volatilité sur les deux indices (pour 100 de notionnel) et vendeur de 0.143 et 0.279 de delta pour 1~\%, que le desk couvre en achetant les indices.
\end{example}

\section{Réserves et bid-offer de paramètre}

La valeur comptabilisée est un cours médian. Le prix auquel la firme pourrait se dégager est moins favorable, pour trois raisons, chacune donnant lieu à une réserve.

\begin{definition}[Réserve de valorisation]\label{def:dv:managing-an-exotic-book:reserve}
Une \emph{réserve de valorisation}\index{réserve de valorisation} est un montant dont une firme diminue la valeur comptable d'une position en dessous de son cours médian modèle, pour une
source précise d'incertitude sur le prix auquel la position pourrait être clôturée~; les réserves sont calculées position par position, tenues dans les comptes et revues à
chaque date de valorisation.
\end{definition}

\begin{definition}[Réserve bid-offer]\label{def:dv:managing-an-exotic-book:bidoffer}
Une \emph{réserve bid-offer}\index{réserve bid-offer} est le coût de clôture du risque net d'un livre au cours acheteur ou au cours vendeur du marché plutôt qu'au cours médian~: l'exposition nette
dans chaque compartiment de risque, multipliée par la demi-fourchette bid-offer des instruments qui serviraient à la clôturer.
\end{definition}

\begin{definition}[Réserve de modèle]\label{def:dv:managing-an-exotic-book:model}
Une \emph{réserve de modèle}\index{réserve de modèle} couvre l'incertitude liée au choix du modèle~: la valeur comptabilisée avec le modèle retenu, diminuée d'un point prudent parmi les valeurs
que prend la position selon un ensemble de modèles alternatifs utilisés par les acteurs du marché.
\end{definition}

\begin{definition}[Bid-offer de paramètre]\label{def:dv:managing-an-exotic-book:param}
Le \emph{bid-offer de paramètre}\index{bid-offer de paramètre} d'une position est la réserve constituée pour une donnée non observable~: la valeur comptabilisée diminuée d'un point prudent parmi
les valeurs prises par la donnée sur sa plage plausible.
\end{definition}

La norme européenne de valorisation prudente fixe le point prudent pour les fonds propres réglementaires~: lorsqu'une plage de valeurs plausibles existe, c'est le point où la firme est
confiante à 90~\% de pouvoir sortir à cette valeur ou à une meilleure. La construction applique la même règle aux trois réserves.

\begin{example}[Trois réserves sur le titre]\label{ex:dv:managing-an-exotic-book:reserves}
\emph{Bid-offer.} La clôture du \omterm{def:dv:greeks-and-the-hedging-pnl:greeks}{vega} à un demi-point de fourchette et celle du delta à 2 points de base coûtent 0.266 et 0.346 pour les deux \omterm{def:dv:greeks-and-the-hedging-pnl:greeks}{vegas} et moins de 0.01 pour les deltas~: soit 0.621 au total
pour 100. \emph{Modèle.} Quatre modèles valorisent le titre à 98.00 (volatilités plates à la monnaie, modèle retenu), 94.33 (volatilité plate au strike de knock-in, trois points
de plus), 98.62 (volatilité plate à celui du déclencheur, un demi-point de moins) et 97.98 (volatilité incertaine, $\pm20\%$). Le point prudent est le plus élevé de ces valeurs pour une
position vendeuse, soit 98.62~: une réserve de 0.621. \emph{Corrélation.} Sur une plage plausible de 0.35 à 0.65, le titre vaut entre 97.64 et 98.50, et le point à 90~\% est de 98.43~: une
réserve de 0.433 (\cref{fig:dv:managing-an-exotic-book:rho,fig:dv:managing-an-exotic-book:models}).
\end{example}

\begin{omfigure}
\begin{tikzpicture}
\begin{axis}[width=10.5cm,height=5cm,xlabel={corrélation entre les deux indices},ylabel={valeur du titre (pour 100)},
  tick label style={font=\small},label style={font=\small},
  xmin=0.18,xmax=0.82,ymin=97.2,ymax=99.2,grid=major,grid style={gray!25},
  table/col sep=comma]
\fill[omProp!15] (axis cs:0.35,97.2) rectangle (axis cs:0.65,99.2);
\fill[omThm!30] (axis cs:0.47,97.2) rectangle (axis cs:0.53,99.2);
\addplot[omThm,very thick,mark=*,mark size=1.4pt] table[x=rho,y=value]{figdata/derivatives/27-managing-an-exotic-book/rho.csv};
\addplot[black,dashed] coordinates {(0.18,98.433) (0.82,98.433)};
\node[font=\footnotesize,anchor=south west] at (axis cs:0.19,98.44) {point à 90~\% à l'origine};
\end{axis}
\end{tikzpicture}

\omcaption{Valeur du titre worst-of à trois ans en fonction de la corrélation de ses deux indices~: la plage plausible à l'origine (claire, de 0.35 à 0.65) et une fois que la
corrélation est observable (foncée, de 0.47 à 0.53). La firme est vendeuse du titre, si bien qu'une corrélation plus élevée lui est défavorable. Données~: le
tutoriel.}\label{fig:dv:managing-an-exotic-book:rho}
\end{omfigure}

\begin{omfigure}
\begin{tikzpicture}
\begin{axis}[width=10.5cm,height=5cm,xbar,bar width=10pt,xlabel={valeur du titre (pour 100)},
  tick label style={font=\small},label style={font=\small},
  xmin=93,xmax=99.5,ymin=-0.6,ymax=3.6,ytick={0,1,2,3},
  yticklabels={volatilité plate à la monnaie (retenu),volatilité plate au strike de knock-in,volatilité plate au déclencheur,volatilité incertaine},
  grid=major,grid style={gray!25},table/col sep=comma]
\addplot[omDef,fill=omDef!40] table[x=value,y=i]{figdata/derivatives/27-managing-an-exotic-book/models.csv};
\draw[black,thick,dashed] (axis cs:98.621,-0.6) -- (axis cs:98.621,3.6);
\end{axis}
\end{tikzpicture}

\omcaption{Le titre sous quatre modèles de volatilité. Pour la position vendeuse, le point prudent (en pointillés) est la valeur la plus élevée, soit un demi-point de volatilité de \omterm{def:dv:greeks-and-the-hedging-pnl:greeks}{vega} au-dessus du
modèle retenu. Données~: le tutoriel.}\label{fig:dv:managing-an-exotic-book:models}
\end{omfigure}

Les réserves ne sont pas simplement additionnées pour le calcul du capital. La première méthode d'agrégation de la norme européenne compte pour moitié, pour chaque position, les ajustements de prix de marché, de clôture et de risque de
modèle, au titre de la diversification entre positions~: 0.838 ici.

\section{Risque de modèle sur le desk}

Le livre~6 traite le risque de modèle comme un processus de gouvernance. Sur le desk, il revêt trois visages. Le premier est le choix entre des modèles qui s'accordent tous sur les vanilles et divergent sur
l'exotique, la plage de la \omterm{def:dv:managing-an-exotic-book:model}{réserve de modèle}~: ici près de quatre points de valeur du titre entre les modèles plats à deux strikes. Le deuxième est un modèle utilisé hors de son
domaine, comme un pricer à volatilité plate pour un knock-in dont la valeur tient au skew. Le troisième est une calibration qui fait varier la valeur sans mouvement de marché
(\omterm{def:dv:fourier-pricing-and-calibration-engineering:recalib}{P\&L de recalibration} du chapitre 24).

\begin{dated}{2026-09}{Guide sur le risque de modèle aux États-Unis}\label{dat:dv:managing-an-exotic-book:mrm}
Le 17 avril 2026, la Federal Reserve, l'OCC et la FDIC ont remplacé le guide de 2011 sur le risque de modèle (SR 11-7) par un guide révisé (lettre SR 26-2 de la Federal Reserve).
Il définit un modèle comme «~une méthode, un système ou une approche quantitative complexe qui applique des théories statistiques, économiques ou financières pour transformer des données d'entrée en
estimations quantitatives~», et laisse hors de son champ les modèles d'IA générative et agentique.
\end{dated}

Un desk vit avec le risque de modèle en valorisant chaque exotique selon au moins deux modèles, en constituant une réserve pour l'écart et en couvrant ce sur quoi les modèles s'accordent. Lorsque les modèles
divergent sur la couverture (le delta d'un worst-of sous deux hypothèses de corrélation, le \omterm{def:dv:greeks-and-the-hedging-pnl:greeks}{vega} d'une barrière sous \omterm{def:dv:local-volatility:model}{volatilité locale} et volatilité stochastique), le desk couvre
la part sur laquelle ils s'accordent et limite le reste.

\section{Stress et concentration}

Les \omterm{def:dv:greeks-and-the-hedging-pnl:greeks}{grecques} décrivent les petits mouvements. Un portefeuille d'\omterm{def:dv:autocallables:autocallable}{autocallables} vendus porte son risque dans les grands~: les barrières de knock-in se regroupent, et les couvertures justes à l'origine
basculent à mesure que les indices s'en approchent.

\begin{example}[Une grille de stress]\label{ex:dv:managing-an-exotic-book:stress}
Le P\&L de la société sur le titre vendu, couvert en delta, en millions pour 100 millions de notionnel, sous des mouvements conjoints des deux indices et des deux volatilités~:
\begin{center}\small
\begin{tabular}{lrrrrr}
\toprule
& $-30\%$ & $-20\%$ & $-10\%$ & $0$ & $+10\%$\\
\midrule
volatilités $-5$ points & 7.02 & $-2.27$ & $-6.31$ & $-5.89$ & $-2.85$\\
inchangées & 13.09 & 5.63 & 1.31 & 0.00 & 0.96\\
volatilités $+5$ points & 18.21 & 12.15 & 8.08 & 6.08 & 5.83\\
\bottomrule
\end{tabular}
\end{center}
La société est acheteuse du put knock-in des investisseurs, si bien qu'un krach accompagné d'une hausse de la volatilité est un gain. La perte vient d'une baisse modérée s'accompagnant d'un recul de la volatilité~: $-6.31$ millions à
$-10\%$ et cinq points de volatilité en moins (\cref{fig:dv:managing-an-exotic-book:stress}).
\end{example}

\begin{omfigure}
\begin{tikzpicture}
\begin{axis}[width=10.5cm,height=5cm,xlabel={mouvement des deux indices (\%)},ylabel={P\&L (millions)},
  tick label style={font=\small},label style={font=\small},
  xmin=-32,xmax=12,ymin=-8,ymax=20,grid=major,grid style={gray!25},
  legend columns=3,legend cell align=left,legend style={font=\footnotesize,at={(0.5,-0.3)},anchor=north,draw=none,fill=none,/tikz/every even column/.append style={column sep=8pt}},
  table/col sep=comma]
\addplot[omProp,very thick,mark=*,mark size=1.4pt] table[x=shock,y=v-5]{figdata/derivatives/27-managing-an-exotic-book/stress.csv};
\addplot[black,thick,mark=square*,mark size=1.4pt] table[x=shock,y=v+0]{figdata/derivatives/27-managing-an-exotic-book/stress.csv};
\addplot[omThm,very thick,mark=triangle*,mark size=1.6pt] table[x=shock,y=v+5]{figdata/derivatives/27-managing-an-exotic-book/stress.csv};
\legend{volatilités $-5$ points,inchangées,$+5$ points}
\end{axis}
\end{tikzpicture}

\omcaption{Le titre worst-of vendu, couvert en delta, sous des mouvements conjoints des indices et de leurs volatilités (100 millions de notionnel). Le pire cas du portefeuille est une baisse modérée
s'accompagnant d'un recul de la volatilité, et non un krach. Données~: le tutoriel.}\label{fig:dv:managing-an-exotic-book:stress}
\end{omfigure}

La concentration est l'autre visage de la taille. Le \omterm{def:dv:greeks-and-the-hedging-pnl:greeks}{vega} du portefeuille sur le second indice, 693\,000 par point, est une position importante face à la volatilité de long terme qui se négocie.
Avec une hypothèse de 100\,000 par point et par jour et une participation d'un dixième, la liquidation prendrait 69 jours. La norme de valorisation prudente demande un
ajustement de concentration lorsque la période de sortie prudente dépasse dix jours. Quand de nombreuses banques vendent les mêmes titres, la même exposition figure dans chaque portefeuille, et la sortie devient encombrée. Les pertes
sur de tels portefeuilles sont publiques lorsqu'elles sont importantes. Natixis a publié un impact non récurrent de 259~millions d'euros sur les dérivés actions asiatiques au quatrième trimestre 2018,
après avoir identifié comme déficiente une stratégie de couverture en Asie.

\section{P\&L du premier jour et sa reprise}

\begin{definition}[P\&L du premier jour]\label{def:dv:managing-an-exotic-book:dayone}
Le \emph{P\&L du premier jour}\index{P\&L du premier jour} d'une transaction est l'écart, à la conclusion, entre son prix de transaction et sa fair value~; selon la norme IFRS 9, il est
comptabilisé immédiatement seulement si la fair value est attestée par un prix de niveau 1 ou par une valorisation n'utilisant que des données observables, et il est sinon différé et comptabilisé
ultérieurement, seulement à mesure que les facteurs qu'intégreraient les intervenants de marché dans leur prix, le temps compris, évoluent.
\end{definition}

La règle ne dit pas comment répartir une marge dont la valorisation utilise des données à la fois observables et non observables. La politique de la construction n'en est qu'une lecture. La \omterm{def:dv:managing-an-exotic-book:bidoffer}{réserve bid-offer}
sur données observables, fait partie de la fair value dès le premier jour. Les réserves sur données non observables, de modèle et de corrélation, sont différées. Ce qui reste de la
marge après ces deux réserves est comptabilisé. Une politique plus stricte diffère la marge tout entière jusqu'à ce que la dernière donnée non observable puisse être observée.

\begin{example}[Le profit différé]\label{ex:dv:managing-an-exotic-book:dayone}
Pour 100 millions de titres vendus au pair~: marge 2.000 millions~; \omterm{def:dv:managing-an-exotic-book:bidoffer}{réserve bid-offer} 0.621~; différé 1.054 (modèle 0.621, corrélation 0.433)~; comptabilisé à la conclusion
0.325. Après un an, les deux indices à 95~\% et le titre toujours en cours, la corrélation se négocie dans une fourchette de 0.47--0.53, et il reste deux ans au titre. La réserve
de corrélation tombe à 0.071 et la \omterm{def:dv:managing-an-exotic-book:model}{réserve de modèle} à 0.559~: 0.424 millions sont repris, dont 0.362 issus de la corrélation. Après deux ans, avec les indices toujours à 95~\%,
un autre 0.336 est repris, et les 0.295 restants à l'échéance (\cref{fig:dv:managing-an-exotic-book:release}).
\end{example}

\begin{omfigure}
\begin{tikzpicture}
\begin{axis}[width=10.5cm,height=5cm,ybar,bar width=14pt,xlabel={année},ylabel={millions},
  tick label style={font=\small},label style={font=\small},
  xmin=-0.6,xmax=3.6,ymin=0,ymax=1.2,xtick={0,1,2,3},grid=major,grid style={gray!25},
  legend columns=2,legend cell align=left,legend style={font=\footnotesize,at={(0.5,-0.3)},anchor=north,draw=none,fill=none,/tikz/every even column/.append style={column sep=8pt}},
  table/col sep=comma]
\addplot[omDef,fill=omDef!40] table[x=year,y=deferred]{figdata/derivatives/27-managing-an-exotic-book/release.csv};
\addplot[omThm,fill=omThm!40] table[x=year,y=released]{figdata/derivatives/27-managing-an-exotic-book/release.csv};
\legend{solde différé,repris dans l'année}
\end{axis}
\end{tikzpicture}

\omcaption{La part différée de la marge du premier jour du titre~: 1.054 millions mis de côté à la conclusion, repris à mesure que la corrélation devient observable et que le titre se rapproche de son échéance
(indices supposés à 95~\% à chaque anniversaire). Données~: le tutoriel.}\label{fig:dv:managing-an-exotic-book:release}
\end{omfigure}

Le report de comptabilisation n'est pas une perte. Il déplace le profit de l'année de vente du titre vers les années où son incertitude est levée. Il modifie pourtant les comportements~: un desk rémunéré sur
le P\&L comptabilisé est moins enclin à conclure des transactions dont la marge repose sur des données non observables.

\section{Tutoriel~: constituer les réserves d'un autocallable}\label{tut:dv:managing-an-exotic-book:tut}

\begin{tutorial}
\textbf{Objectif.} Valoriser un worst-of \omterm{def:dv:autocallables:autocallable}{autocallable}, fixer son coupon pour une marge de 2~\%, calculer ses \omterm{def:dv:managing-an-exotic-book:bidoffer}{réserves bid-offer}, de modèle et de corrélation, répartir son \omterm{def:dv:managing-an-exotic-book:dayone}{P\&L du premier jour} et
programmer sa reprise, et lancer une grille de stress. \textbf{État final~:} les quatre chiffres, le tableau de stress et les résultats du problème du week-end.
\begin{tutsteps}
\item \textbf{Le point prudent} et les \omterm{def:dv:managing-an-exotic-book:model}{réserves de modèle} et de paramètre~:
\omcode{code/firm/reserves/firm_reserves.py}{22}{47}{Point prudent, réserve de modèle et bid-offer de paramètre.}\label{lst:dv:managing-an-exotic-book:prudent}
\item \textbf{La répartition du premier jour} et l'échéancier de reprise~:
\omcode{code/firm/reserves/firm_reserves.py}{56}{81}{P\&L du premier jour et sa reprise.}\label{lst:dv:managing-an-exotic-book:dayone}
\item \textbf{Lancer} \texttt{dv\_reserves.fair\_coupon()}, \texttt{reserves\_at(0)}, \texttt{day\_one\_study()}, \texttt{rho\_curve()}, \texttt{stress()} et
\texttt{fig\_reserves.py}.
\end{tutsteps}
\textbf{Pour aller plus loin.} Ajouter le pricer à \omterm{def:dv:local-volatility:model}{volatilité locale} du chapitre 18 à l'ensemble de modèles~; valoriser la reprise avec les indices à 80~\% après un an~; fixer la
plage de corrélation à partir de la dispersion des \omterm{def:dv:multi-asset-options:implied}{corrélations implicites} du chapitre 17.
\end{tutorial}

\section{Construction~: le calculateur de réserves}\label{bld:dv:managing-an-exotic-book:reserves}

\begin{build}
\bfield{Rôle} Le calculateur de réserves de la firme miniature~: bid-offer par compartiment de risque, \omterm{def:dv:managing-an-exotic-book:model}{réserves de modèle} à partir d'un ensemble de modèles, \omterm{def:dv:managing-an-exotic-book:param}{bid-offer de paramètre} sur données non observables,
agrégation, répartition du premier jour et sa reprise, grilles de stress et concentration.
\bfield{Interface} \texttt{bid\_offer\_reserve(exposures, half\_spreads)}~; \texttt{prudent\_point(values, confidence, weights)}~; \texttt{model\_reserve(values\_by\_model,
booked)}~; \texttt{parameter\_reserve(value\_of, mid, lo, hi)}~; \texttt{aggregate(reserves, diversification)}~; \texttt{day\_one(price, booked\_liability, observable,
unobservable)} $\to$ \texttt{DayOne}~; \texttt{release\_schedule}~; \texttt{stress\_grid(value\_of, spot\_shocks, vol\_shocks, hedge\_delta)}~; \texttt{concentration\_days}.
\bfield{Règles} Les valeurs sont la valeur du portefeuille pour la société, et les réserves sont positives~; le point prudent est le point à 90~\% de la plage plausible~; les réserves non observables sont
différées, les réserves observables font partie de la fair value~; chaque réserve est recalculée à chaque date de valorisation.
\bfield{Tests d'acceptation} \texttt{code/firm/reserves/tests/}~: le point prudent d'une plage uniforme~; une \omterm{def:dv:managing-an-exotic-book:model}{réserve de modèle} nulle lorsque le modèle retenu est déjà le modèle prudent~;
la réserve de paramètre d'une valeur linéaire sur une plage uniforme~; bid-offer, agrégation et arithmétique du premier jour~; un échéancier de reprise dont la somme égale le différé~; une grille de stress sur un
portefeuille linéaire.
\bfield{Pour aller plus loin} Réserves sur un portefeuille avec compensation des expositions opposées~; un ajustement de concentration à partir de la période de sortie~; un ensemble de modèles issu des chapitres 9, 10 et des
pricers du chapitre 18, via la bibliothèque du chapitre 28.
\end{build}

\begin{omsources}
\item Commission Regulation (EU) No 1255/2012 (IFRS 13 \emph{Fair Value Measurement}), paragraphes 76--86.
\item Commission Regulation (EU) 2016/2067 (IFRS 9 \emph{Financial Instruments}), paragraphe B5.1.2A.
\item Commission Delegated Regulation (EU) 2016/101 (valorisation prudente), articles 9, 11 et 14 et annexe.
\item Board of Governors of the Federal Reserve System, SR 26-2, «~Revised guidance on model risk management~» (17 avril 2026)~; OCC Bulletin 2026-13.
\item Natixis, 2018 fourth-quarter and annual results, communiqué de presse, 12 février 2019.
\end{omsources}

\section{Exercices}

\begin{exercise}[$\star$]\label{exo:dv:managing-an-exotic-book:1}
Expliquer pourquoi la note worst-of est un instrument de niveau 3, et dans quel cas elle pourrait passer au niveau 2.
\end{exercise}

\begin{exercise}[$\star$]\label{exo:dv:managing-an-exotic-book:2}
Expliquer pourquoi le point prudent est la valeur de modèle la plus élevée pour la société, qui est vendeuse de la note.
\end{exercise}

\begin{exercise}[$\star$]\label{exo:dv:managing-an-exotic-book:3}
Vérifier que la répartition du \omterm{def:dv:managing-an-exotic-book:dayone}{P\&L du premier jour} égale bien la marge, et la calculer selon une politique qui diffère toute la marge nette de bid-offer.
\end{exercise}

\begin{exercise}[$\star\star$]\label{exo:dv:managing-an-exotic-book:4}
La \omterm{def:dv:managing-an-exotic-book:model}{réserve de modèle} (0.621) est égale au bid-offer de \omterm{def:dv:greeks-and-the-hedging-pnl:greeks}{vega} (0.621) à trois décimales près. Expliquer s'il y a double comptage.
\end{exercise}

\begin{exercise}[$\star\star$]\label{exo:dv:managing-an-exotic-book:5}
Expliquer pourquoi une corrélation plus élevée augmente la valeur d'une note worst-of.
\end{exercise}

\begin{exercise}[$\star\star$]\label{exo:dv:managing-an-exotic-book:6}
Expliquer pourquoi la case la plus défavorable de la grille de stress correspond à une baisse modérée accompagnée d'une baisse de la volatilité.
\end{exercise}

\begin{exercise}[$\star\star\star$]\label{exo:dv:managing-an-exotic-book:7}
\emph{Programmation.} Calculer la réserve de corrélation avec la plage plausible 0.2--0.8 au lieu de 0.35--0.65. Déterminer ce que devient le \omterm{def:dv:managing-an-exotic-book:dayone}{P\&L du premier jour}.
\end{exercise}

\begin{exercise}[$\star\star\star$]\label{exo:dv:managing-an-exotic-book:8}
\emph{Repérer l'erreur.} «~Nos réserves sont prudentes~: nous les additionnons toutes pour chaque position, sans aucun bénéfice d'agrégation, de sorte que les fonds propres ne sont jamais sous-estimés.~»
\end{exercise}

\section{Problème~: Le profit différé}

\begin{problem}[{Problème du week-end --- quelle part de la marge est un profit}]\label{pb:dv:managing-an-exotic-book:1}
Le desk de l'ouverture vend 100 millions de l'\omterm{def:dv:autocallables:autocallable}{autocallable} worst-of à trois ans au pair~; son modèle valorise les notes à 98.

\medskip\noindent\textbf{Partie I --- La note.}
\begin{enumerate}
\item Donner le coupon qui donne à la note une valeur de 98, et dire quelles sont les positions longues et courtes de la société une fois la note vendue.
\item Indiquer quelles données de la valorisation sont observables et lesquelles ne le sont pas.
\item Indiquer à quel niveau de la hiérarchie de la fair value se situe la note, et ce que dit IFRS 9 sur son gain du premier jour.
\item Donner les \omterm{def:dv:greeks-and-the-hedging-pnl:greeks}{vegas} et les deltas de la position de la société.
\item Indiquer ce que la société couvre et ce qu'elle ne peut pas couvrir.
\end{enumerate}

\medskip\noindent\textbf{Partie II --- Réserves.}
\begin{enumerate}[resume]
\item Calculer la \omterm{def:dv:managing-an-exotic-book:bidoffer}{réserve bid-offer}.
\item Donner les valeurs des quatre modèles et la \omterm{def:dv:managing-an-exotic-book:model}{réserve de modèle}.
\item Donner l'intervalle de valeurs de la note sur la plage plausible de la corrélation, ainsi que la réserve de corrélation.
\item Indiquer le total agrégé selon la première méthode de la norme européenne.
\item Déterminer combien de temps il faudrait pour solder le \omterm{def:dv:greeks-and-the-hedging-pnl:greeks}{vega} du deuxième indice, et ce qui en découle.
\end{enumerate}

\medskip\noindent\textbf{Partie III --- Premier jour et après.}
\begin{enumerate}[resume]
\item Répartir la marge entre la part facturée, la part différée et la part reconnue.
\item Déterminer ce qui change après un an, et ce qui est repris.
\item Déterminer la part de la reprise provenant de la corrélation.
\item Déterminer l'effet sur la reprise d'une baisse des indices de 80~\% la première année.
\item Indiquer comment la case la plus défavorable de la grille de stress modifierait la couverture du desk.
\end{enumerate}

\medskip\noindent\textbf{Partie IV --- Jugement.}
\begin{enumerate}[resume]
\item Dire si différer la marge est prudent ou simplement lent.
\item Expliquer comment le report modifie les incitations du desk.
\item Indiquer qui doit assumer la responsabilité de l'ensemble de modèles et des plages plausibles.
\item Énoncer le \emph{résultat nommé}~: le \omterm{def:dv:managing-an-exotic-book:dayone}{P\&L du premier jour} reconnu à la conclusion sur un \omterm{def:dv:autocallables:autocallable}{autocallable} worst-of à trois ans, et le montant repris après un an, lorsque
la donnée de corrélation devient observable.
\item Dire en une phrase à quoi sert une réserve.
\end{enumerate}
\end{problem}

\section{Questions d'entretien}

\begin{interviewq}[$\star$ \iqroles{banque, risque}]\label{iq:dv:managing-an-exotic-book:1}
Qu'est-ce que le \omterm{def:dv:managing-an-exotic-book:dayone}{P\&L du premier jour}, et quand une banque peut-elle le reconnaître~?
\end{interviewq}

\begin{interviewq}[$\star\star$ \iqroles{risque}]\label{iq:dv:managing-an-exotic-book:2}
Comment calculeriez-vous une \omterm{def:dv:managing-an-exotic-book:model}{réserve de modèle} pour un produit exotique~?
\end{interviewq}

\begin{interviewq}[$\star\star$ \iqroles{banque, trader}]\label{iq:dv:managing-an-exotic-book:3}
Quelles expositions une banque qui émet des \omterm{def:dv:autocallables:autocallable}{autocallables} conserve-t-elle, et pourquoi sont-elles difficiles à couvrir~?
\end{interviewq}

\begin{interviewq}[$\star\star$ \iqroles{risque}]\label{iq:dv:managing-an-exotic-book:4}
Concevriez-vous un test de stress pour un portefeuille d'\omterm{def:dv:autocallables:autocallable}{autocallables} vendus à découvert~? Quel scénario redoutez-vous le plus~?
\end{interviewq}

\begin{interviewq}[$\star\star$ \iqroles{développeur, risque}]\label{iq:dv:managing-an-exotic-book:5}
Que doit conserver un système de calcul des réserves pour qu'un auditeur puisse reproduire une réserve un an plus tard~?
\end{interviewq}

\begin{interviewq}[$\star\star\star$ \iqroles{banque}]\label{iq:dv:managing-an-exotic-book:6}
Une corrélation non observable devient observable. Que devient le P\&L du portefeuille, et pourquoi~?
\end{interviewq}
